Operational deployment imposes practical requirements on converting an anomaly score stream into a sequence of binary anomaly labels. When large streams are monitored continuously, even a low false-positive rate can produce a substantial investigation workload \cite{msl_smap_data}. A single incident may also generate intermittent score excursions, leading to repeated alarm episodes under pointwise thresholding. Suppressing short-lived fluctuations can reduce nuisance alerts but may also delay detection or obscure brief incidents.

A practical decision rule must therefore balance prompt detection, false-alarm burden, and the consolidation of related evidence into coherent alarm episodes. To preserve the unsupervised nature of the pipeline, this trade-off must be calibrated without optimizing against observed detection delays or missed-detection rates, whose direct measurement requires labeled anomalies. Additionally, this removes the need to collect and manually annotate anomalous examples for calibration \cite{expensivelabeling}. An open challenge remains: developing decision rules for streaming anomaly detection that satisfy these operational requirements using normal data alone, without assuming prior knowledge of the anomalies to be encountered.
